\summarychapter{17}{God's Consul}{Gregory the Great}

\begin{inbrief}
Europe rose from the ruins of the empire, and more than any other force Christianity brought order
out of the chaos. Its instruments were the missionary monks, the papal office and the theology of
Augustine, and no one mastered all three better than Pope Gregory I. The chapter follows his life, his
practical rule of Rome and Italy, his mission to England and the popular faith he handed on to the
Middle Ages.
\question{What did Christianity bring to the devastation to build the new order called Christian
Europe?}
\end{inbrief}

Early in 590 Rome, already struck by floods and war, was ravaged by plague. Carts were piled with the
dead, and Pope Pelagius II himself died in agony. For six months there was no pope. The clergy then
elected the monk Gregory, who fled and hid in the forest until he was dragged back and consecrated on
September 3, 590. Fifty, balding and frail, he began with a public act of penitence, three days of
processions through the streets. The plague at last subsided, and legend told that the archangel
Michael sheathed his sword above Hadrian's mausoleum, ever since called the Castle of St.~Angelo.

\ssection{Gregory's Place in History}
The church enlisted Celtic and Benedictine monks as a spiritual militia, turned to the papacy for a
stable structure and kept Augustine's ideals as its framework of meaning. Gregory's \emph{Pastoral
Care} teaches that a spiritual leader must neither lose the inner life in outward cares nor neglect
outward duties for the inner life: Christ prayed on the mountain and worked miracles in the cities.
The words describe their author. When he died in 604 his epitaph called him ``God's Consul'': wholly
God's, yet ruling church and world like a Roman statesman. He was soon called ``the Great'' and ranked
with Augustine, Ambrose and Jerome as a Latin Father. In intellect he hardly belongs with them, but he
joined great executive ability to warm sympathy, and no one represents the early Middle Ages better.

Born about 540 into an old senatorial family and trained for government, Gregory grew up while Rome
changed hands repeatedly. Narses became Justinian's viceroy of Italy in 554; after a few years of
peace the Lombards began burning churches, killing bishops and laying the fields waste. The city of
the caesars was becoming the city of the popes. At thirty-three, in 573, he was appointed prefect of
Rome, responsible for grain, poor relief and public works. Uneasy with worldly power, he resigned
after a few years. On his father's death he founded seven monasteries, gave the rest of his fortune
to the poor, turned the family palace into the monastery of St.~Andrew and lived as a monk with such
austerity that he ruined his health; yet he called those years his happiest. In 579 Pelagius II made
him one of the seven deacons and sent him as ambassador to Constantinople. He returned in 585 and
became abbot of St.~Andrew, until the plague carried off the pope.

\ssection{Picking Up the Pieces of a Shattered World}
Soon after his election the Lombards besieged Rome, and Gregory broke off his sermons on Ezekiel in
grief. Like other serious men he believed the end of the world was near: cities destroyed, fields
waste, people carried off or killed. ``If we love such a world, we love not our joys, but our
wounds.'' The Roman church was almost the only organized institution in the West to survive. Though he
regarded his election as a punishment, Gregory threw himself into the work. He wrote sharp, practical
letters to the managers of his estates and to his bishops and set out the principles of ministry in
his \emph{Pastoral Rule}. Nothing was too great or too small for his attention, though gout often
kept him in bed: ``I am daily dying, but never die.'' The prestige of the medieval papacy rests in
large part on this practical government.

The Roman church owned wide lands around Rome, in southern Italy and in Sicily, the ``patrimony of
St.~Peter,'' some 1,800 square miles, which made it the richest landowner in Italy. When the Lombards
destroyed the imperial administration, the officials of the patrimony fed the people and collected
the taxes, with Gregory at their head. He also organized the defense of central Italy, appointed a
military governor and made peace with two Lombard leaders, so that after 595 the pope counted for
more in Lombard politics than any imperial officer. From then on the pope was a political figure as
well as a Christian leader.

His view of the office appears in his clash with John IV, patriarch of Constantinople, who used the
honorary title ``universal bishop,'' confirmed by a synod of 588. Gregory denounced it as proud,
profane and diabolical, threatened to break communion and appealed to the emperor. For himself he
chose ``servant of the servants of God,'' which became a standing papal title, and he refused to be
called universal pope. Yet he claimed and exercised, so far as he could, oversight of the whole
church, and his successors took prouder titles without scruple. Pride haunted him. His \emph{Moralia}
on Job calls it the root of the vices, a disease that corrupts the whole body, and Shelley suggests
that his restless activity sprang from the struggle against his own pride, a struggle in keeping with
his monastic calling.

\ssection{Missionary Statesman}
With Gregory monasticism first ascended the papal throne. He kept a monk's simplicity, surrounded
himself with monks, made them bishops and legates, confirmed Benedict's Rule at a Roman council and
protected the liberty and property of the monasteries. A story, part fact and part fiction, tells how
before he was pope he saw English boys for sale in the Roman market. Told that they were Angles, he
said that they looked like angels and should be heirs with the angels; that coming from Deira they
should be saved from God's wrath (\emph{dei ira}); and that in the land of King Aella Alleluia must be
sung. Prevented from going himself, he later sent the Benedictine Augustine with forty monks to
England. The success of that mission in Kent links all Anglo-American Christianity with the early
church.

\ssection{Defender of Orthodoxy}
Gregory's teachers were Ambrose, Augustine and Jerome. He added no new ideas, but he formulated the
common faith of his day, including the crude and sometimes pagan notions of the unlettered, lent it
his authority and so made it part of the faith of the Western church. Medieval thought cannot be
followed without him. Adam's fall, he taught, weakened human freedom without destroying it. Moved by
grace, a person can cooperate with it and gain merit by good works, the joint product of grace and
will. Baptism forgives freely, but sins after baptism must be atoned for by penance, in which the
sinner punishes himself rather than leave it to God. Penance requires sincere repentance, confession
and meritorious works involving sacrifice, such as alms, ascetic practices and prayer, and whether
one has done enough cannot be known before death.

Sinners have help. Gregory made central the older belief in the intercession of the saints, who are
ready to be advocates before the severe judge Jesus and ``entreat that they may be entreated.'' He
encouraged the collecting and veneration of relics---hair, fingernails, pieces of clothing---believed to
have great powers. What is not atoned for in this life can be purged in purgatory, a place
of purification and suffering for those who die neither perfectly holy nor wicked but with minor sins
still to pay for. The supreme miracle is the Eucharist. Christ's body and blood are really present in
the bread and wine, and the Eucharist is a sacrifice offered by the priest for the sins of those who
take part or for whom it is offered. It works like penance, replacing suffering otherwise due, and
helps the dead in purgatory as well as the living. In his \emph{Dialogues} Gregory tells of Justus, a
monk who died punished for hoarding money. Gregory ordered Mass said for him on thirty consecutive
days; on the thirtieth Justus appeared to his brother Copiosus, who knew nothing of the Masses, and
said that he had that morning been admitted to communion. The doctrine was widely accepted in the West
from then on and gave medieval Christianity its particular tone.

\begin{keyfacts}
\keypeople{Gregory I (the Great); Pope Pelagius II; Narses; the Lombards; John IV, patriarch of
Constantinople; Augustine of Canterbury.}
\keydates{%
  \kd{c. 540}{Birth of Gregory the Great in a senatorial family of Rome}
  \kd{573}{Gregory, the future pope, prefect of Rome}
  \kd{579}{Gregory sent to Constantinople as papal ambassador}
  \kd{590}{Gregory the Great consecrated pope (September 3)}
  \kd{595}{Gregory makes peace with the Lombards; the pope a political force in Italy}
  \kd{604}{Death of Gregory the Great, ``God's Consul''}}
\keyterms{Patrimony of St.~Peter; ``servant of the servants of God''; universal bishop; Latin Fathers;
\emph{Pastoral Care}; merit; penance; intercession of the saints; relics; purgatory; sacrifice of the
Mass.}
\end{keyfacts}
